\chapter{Future Work}

\section{A structural account of $A[\sqrt{n}]$}

The development of $A[\sqrt{n}]$ in this document is deliberately elementary:
elements are pairs, and the ring and field laws are checked coordinatewise. That
is easy to follow and easy to compute with, but it obscures a good deal of
structure, and the ring proofs feel more laborious than the mathematics
warrants.

What the construction really is, is the \emph{regular representation}. Every
element $z$ of $A[\sqrt{n}]$ acts on $A[\sqrt{n}]$ by multiplication, and that
action is $A$-linear, so in the basis $\{1, \sqrt{n}\}$ it is a matrix over $A$:
\begin{displaymath}
  M(z) \;=\;
  \begin{pmatrix} z_1 & n z_2 \\ z_2 & z_1 \end{pmatrix},
  \qquad z = z_1 + z_2\sqrt{n}.
\end{displaymath}
The second column is where $n$ enters, and it is the only place. Since
$(zw)x = z(wx)$, the map $z \mapsto M(z)$ is a ring homomorphism, and
multiplying the matrices out reproduces the multiplication formula this
development defines by hand. Under this correspondence:
\begin{itemize}
\item the norm $z \cdot \overline{z}$ is $\det M(z)$, so its multiplicativity is
      that of the determinant;
\item $2z_1$ is $\operatorname{tr} M(z)$;
\item conjugation is the adjugate, so the formula for the inverse is Cramer's
      rule, and the fact that the norm vanishes only at zero is the statement
      that a matrix is invertible exactly when its determinant is not;
\item the characteristic polynomial is
      $\lambda^2 - 2z_1\lambda + (z_1^2 - n z_2^2)$, whose coefficients lie in
      $A$ while its roots are $z_1 \pm z_2\sqrt{n}$ --- the two conjugate
      embeddings.
\end{itemize}

That last point explains why the ordering can be described at all. The order
asks whether one particular eigenvalue is non-negative, and $\sqrt{n}$ appears
in the roots but never in the coefficients; the criterion of
Definition~\ref{def:posCone} is the standard way of locating a root of a
quadratic from its coefficients. It also explains the shape of the argument for
closure under addition, which still turns on elementary symmetric combinations
of the coordinates through which roots and coefficients communicate. \TODO{Update
the specific pairing here --- the proof now gets the harder half of
closure under addition as a corollary of the easier half, via the $\sqrt n$
rotation, rather than via a standalone multiplier argument.}

This suggests presenting $A[\sqrt{n}]$ as $A[x]/(x^2 - n)$ and inheriting its
structure rather than verifying it. Mathlib's \code{AdjoinRoot} supplies the
commutative ring, the algebra structure, decidable equality, a power basis, and
--- given irreducibility, which is what the non-square hypothesis amounts to ---
the field structure. That would remove essentially all of the ring boilerplate
here, including the one proof in this development that has to navigate
associativity and commutativity by hand. Mathlib also defines \code{Algebra.norm}
as precisely the determinant of the multiplication map, so the norm and its
multiplicativity would come from the same place.

What it would \emph{not} simplify is the ordering, which is the bulk of the work
and was the difficult part. Mathlib supplies the ring and field theory of an
extension but no order theory for one: that is Artin--Schreier territory, and
the results are not there. The matrix picture explains what the sign function is
doing without providing a theorem one can cite in its place.

So the honest assessment is that the refactor is worth doing for the algebra and
irrelevant to the analysis. There is a second reason to want it, though, which
has nothing to do with brevity: the elementary presentation is a completely
concrete instance of a construction usually met in the abstract, and the two
side by side --- pairs of numbers on one hand, a quotient of a polynomial ring
on the other, provably the same --- would make a good introduction to the
commutative algebra involved.

One caveat for whoever attempts it: \code{AdjoinRoot} is a quotient type and its
field structure is noncomputable, so it would serve for reasoning while the
existing pair representation continued to serve for computation, the two related
by a theorem. That is the same division discussed in
Chapter~\ref{chap:arrangements}, and it recurs often enough in this project to
be worth naming as a habit rather than rediscovering each time.

\section{Deferred components}

\TODO{Hardware floating point:
\begin{itemize}
  \item now built on LeanCert dyadic intervals (software, unbounded exponent)
  \item hardware \code{Float}: opaque in Lean; option is IEEE 754 axioms,
    Flocq-style~\cite{flocq}
  \item swap design: see the note after def:filteredInv
  \item performance unmeasured
\end{itemize}}

Extraction of a reasonably efficient executable has not been attempted. The
ordering on $A[\sqrt{n}]$ is computable by construction --- decidability of cone
membership is what makes it so --- but no attention has yet been paid to whether
the compiled code is fast, and some of the algebraic instances inherit
deliberately naive default implementations from Mathlib.

The rendering support is presently limited to producing static diagrams. The
interactive widget infrastructure available in Lean would permit dragging shapes
directly in the editor, which is the natural presentation for the drag-and-drop
application.

% vim: ts=2 sw=2 et ai
